% 中文摘要
\thesisfrontmatterpage{摘要}
\pagenumbering{Roman}

\addcontentsline{toc}{section}{摘要}

\section*{摘\quad 要}
% 摘要：二字间间空两格，黑体三号居中，段前，段后各空一行。
\thesisbodyfont
手内旋转要求灵巧手在不重新抓取物体的情况下持续改变其姿态，是精细操作中的一项基础能力。面向低成本 LEAP Hand 在无视觉、无触觉条件下的连续旋转部署，本文在 Isaac Lab 中建立圆柱体连续手内旋转任务、奖励函数和统一评测协议，并将 HORA/RMA 两阶段适应思路改造为适用于 LEAP 平台的 \(8\) 维技能先验表示。为缓解第二阶段部署策略在分布外圆柱参数下的退化，本文在训练中扩大质心偏移覆盖范围，并引入教师——部署动作一致性损失，最终得到仅依赖本体感觉历史的部署策略。\par
仿真中，第一阶段教师策略能够在 \(20\,\mathrm{s}\) 回合内维持稳定的圆柱体连续旋转。增强后的 \(\mathrm{Deploy\text{-}Refined}\) 在分布内评测中成功率为 \(1.0000\)，在常规分布外评测中将成功率由 \(0.8750\) 提高到 \(0.9141\)，存活时间、净旋转圈数和轴向角速度也同步改善。消融实验显示，动作一致性约束贡献了第二阶段补强的主要增益，放宽质心偏移分布则增加了高风险样本覆盖。实机测试中，同一策略可以在真实 LEAP Hand 上闭环运行，并在标准圆柱体、尺寸变化圆柱体和细高饮料瓶上产生可重复的低速受控旋转。本文结论限于圆柱体参数泛化和少量近圆柱对象迁移，任意多物体稳定旋转仍是后续研究目标。\par
~\\
\hspace*{2em}{\heiti \zihao{-4}关键词：}LEAP Hand；手内旋转；强化学习；本体感觉自适应；Sim2Real\\
% 关键词：宋体12磅，行距20磅，段前段后0磅，关键词之间用分号隔开,关键词三个字加粗。
